\section{Conclusion}
\label{sec:conclusion}

We investigated where rhyming knowledge is encoded in transformer language models through three experiments on \gptmed using 13{,}105 words with IPA transcriptions.
Our layer-wise probing reveals that phonetic information is \emph{front-loaded}: it peaks at the embedding/layer-0 boundary and does not improve through the network, contrasting with semantic and syntactic features that build through deeper layers.
Our geometric analysis shows that each phoneme has a distinct direction in embedding space ($5.3\times$ above random baseline) but these directions are not orthogonal---they form an overlapping phonetic manifold that compresses in deeper layers.
Our heteronym test uncovers a primary-pronunciation bias in \gptfour (92\% primary vs.\ 58\% secondary accuracy), demonstrating that context-dependent phonetic processing remains an open challenge.

The answer to ``where is rhyming?'' is: mostly in the token embedding table, refined by the first transformer layer, and preserved but not enriched through the rest of the network.

Several directions for future work emerge from these findings.
Replicating this analysis on \llama and other model families would test whether the front-loading pattern generalizes.
Causal interventions---modifying phoneme direction vectors and observing changes in rhyming behavior---would establish whether these directions are merely correlated with phonetic content or causally implicated.
Extending the heteronym evaluation to more words and models would better characterize the primary-pronunciation bias and test whether improved prompting or fine-tuning can mitigate it.
Finally, analyzing multi-token words would reveal whether the phonetic encoding patterns we observe for single tokens extend to subword-assembled pronunciations.
